\section{Method}
\label{sec:method}

\subsection{Problem Setting and Notation}

Let patient $i$ contribute an observation set $\mathcal{O}_i=\{(t_{ij},m_{ij},v_{ij})\}_{j=1}^{n_i}$ with observation times $t_{ij}\in[0,T_i]$, variable indices $m_{ij}\in\{1,\dots,d\}$, and values $v_{ij}\in\mathbb{R}$. The set is irregular in the sense that neither the number of observations nor their spacing is controlled by the analyst: a laboratory panel may be drawn four times in an hour and then not for eight hours. A score function $f_\theta$ maps the history available at time $t$, denoted $\mathcal{O}_{i,\le t}$, to a risk score $s_{it}=f_\theta(\mathcal{O}_{i,\le t})\in(0,1)$. The score function is treated as a black box in this section; any of the backbones in \cref{sec:related} may be substituted, and \cref{sec:setup} fixes a small set for the experiments.

The event of interest occurs at time $t^\star_i$, and for patients who never experience it we set $t^\star_i=\infty$. A prediction is useful if it fires in the window $[t^\star_i-W,\,t^\star_i)$ for a lead-time window $W$, which we fix at twelve hours following common practice in deterioration prediction. Because the score is defined at observation times only, all decisions below are indexed by the observation grid of each patient rather than by a global clock.

\subsection{Stateful Alarm Policies}

A threshold rule declares an alarm whenever $s_{it}\ge\lambda$. Such a rule has no memory, and in deployment it fires repeatedly on the same deteriorating patient, which is the mechanism by which alert volume becomes unmanageable. We therefore define an alarm policy $\pi_\lambda=(\lambda,\tau,\kappa)$ with three components: a threshold $\lambda$, a refractory window $\tau$ expressed in observation steps, and an escalation tier $\kappa$ that specifies how many suppressed exceedances must accumulate before the refractory period is overridden. The policy carries an internal state $z_{it}$ and emits
\begin{equation}
\label{eq:policy}
a_{it}=\mathbb{1}\!\left[s_{it}\ge\lambda \ \wedge\ \left(z_{it}=0 \ \vee\ c_{it}\ge\kappa\right)\right],
\qquad
z_{i,t+1}=
\begin{cases}
\tau, & a_{it}=1,\\
\max(0,\,z_{it}-1), & a_{it}=0,
\end{cases}
\end{equation}
where $c_{it}$ counts consecutive exceedances of $\lambda$ that were suppressed. The refractory window is the control that directly trades alarm volume against repeated detection, and the escalation tier prevents a long refractory period from masking a genuinely escalating patient. Setting $\tau=0$ and $\kappa=\infty$ recovers the memoryless threshold rule, so \cref{eq:policy} is a strict generalisation rather than a replacement.

Two properties of \cref{eq:policy} matter for what follows. First, the policy is a deterministic function of the score trajectory given $(\lambda,\tau,\kappa)$, so the only quantity that must be calibrated against deployment data is the scalar $\lambda$; the structural parameters are selected by cross-validation on the source cohort. Second, the emitted sequence $\{a_{it}\}$ is dependent in time by construction, and this dependence is not a nuisance to be averaged away but the mechanism through which the budget is enforced.

\subsection{Why Timestep-Level Calibration Fails}

It is tempting to calibrate $\lambda$ by treating each observation time as an example and applying split conformal prediction to the indicator $a_{it}$. This is invalid for two distinct reasons, and separating them clarifies what the proposed procedure must repair.

The first is granularity. The operational quantity is a rate per patient-day,
\begin{equation}
\label{eq:risk}
R(\lambda)=\mathbb{E}\!\left[\frac{1}{D_i}\sum_{t}a_{it}(\lambda)\right],
\end{equation}
where $D_i$ is the length of patient $i$'s record in days. The expectation is over patients, while the inner sum is over dependent observations. Controlling the marginal alarm probability at a randomly chosen timestep does not control \cref{eq:risk}: the two differ by the inverse of the mean record length and by the within-patient correlation of $a_{it}$.

The second is state dependence. Because $a_{it}$ depends on $a_{i,t-1}$ through $z_{it}$, the timestep-level losses are neither independent nor exchangeable once the policy is active. Calibrating on a pooled timestep sample therefore estimates the alarm rate of a policy evaluated on a different dependence structure from the one that will run in the ward. The error is not conservative in either direction; it depends on how strongly alarms cluster, which is precisely the regime a well-designed refractory window creates.

\subsection{Blocked Conformal Risk Control}

The repair is to move the unit of calibration from the timestep to the patient. Define the patient-level loss
\begin{equation}
\label{eq:blockloss}
L_i(\lambda)=\frac{1}{D_i}\sum_{t}a_{it}(\lambda)\in[0,\,L_{\max}],
\end{equation}
which is exactly the quantity inside the expectation in \cref{eq:risk}. The bound is $L_{\max}=\sup_{H\ge H_{\min}}\frac{24}{H}\left(\left\lfloor\frac{H/\Delta-1}{\tau+1}\right\rfloor+1\right)$, where $H$ is the length of stay, $\Delta$ the sampling interval and $H_{\min}$ the shortest stay admitted to the cohort. The supremum is not attained at one day: with hourly sampling and $\tau=4$ a 26-hour stay permits six alarms, or $144/26\approx5.54$ per day, while a 24-hour stay permits only five. Evaluating the bound at a nominal day therefore understates it, and since short stays are exactly the ones that generate high per-day rates, $H_{\min}$ is a modelling decision that must be declared rather than a detail. Patients are exchangeable draws from the deployment population even though observations within a patient are not, so averaging \cref{eq:blockloss} over patients restores the exchangeability that conformal calibration requires. Let $\Lambda$ be a finite grid of candidate thresholds and let $\widehat{R}_n(\lambda)=\frac{1}{n}\sum_{i=1}^{n}L_i(\lambda)$ be the calibration average over $n$ held-out patients.

\begin{proposition}[Budget validity]
\label{prop:validity}
Fix $\alpha\in(0,1)$ and a budget $B>0$. Let
\begin{equation}
\label{eq:lambda}
\widehat{\lambda}=\inf\left\{\lambda\in\Lambda:\ \widehat{R}_n(\lambda)+L_{\max}\sqrt{\frac{\log(|\Lambda|/\alpha)}{2n}}\ \le\ B\right\},
\end{equation}
and set $\widehat{\lambda}=\max\Lambda$ if the set is empty. If calibration and deployment patients are exchangeable draws from the same population and $L_i$ is bounded as in \cref{eq:blockloss}, then
\begin{equation}
\label{eq:bound}
\Pr\!\left[R(\widehat{\lambda})>B\right]\ \le\ \alpha ,
\end{equation}
where $R(\lambda)=\mathbb{E}[L_i(\lambda)]$ is the population alarm rate defined in \cref{eq:risk}.
\end{proposition}

\begin{proof}
For a fixed $\lambda\in\Lambda$, Hoeffding's inequality applied to the $n$ bounded independent losses gives $\Pr[\widehat{R}_n(\lambda)-R(\lambda)\ge\epsilon]\le\exp(-2n\epsilon^{2}/L_{\max}^{2})$. A union bound over the $|\Lambda|$ candidate thresholds with $\epsilon=L_{\max}\sqrt{\log(|\Lambda|/\alpha)/(2n)}$ yields
\begin{equation}
\Pr\!\left[\exists\lambda\in\Lambda:\ \widehat{R}_n(\lambda)-R(\lambda)\ge\epsilon\right]\le\alpha .
\end{equation}
On the complement event, $\widehat{R}_n(\lambda)\ge R(\lambda)-\epsilon$ simultaneously for every $\lambda$. If $R(\widehat{\lambda})>B$ then $\widehat{R}_n(\widehat{\lambda})>B-\epsilon$, which contradicts the selection rule \eqref{eq:lambda} that requires $\widehat{R}_n(\widehat{\lambda})+\epsilon\le B$. Hence $\{R(\widehat{\lambda})>B\}$ is contained in the complement of the concentration event, whose probability is at most $\alpha$.
\end{proof}

Three remarks follow. The bound is on the population alarm rate, not on the per-patient tail: a ward experiences both, and a guarantee about the mean does not imply that most patients stay within budget when the loss distribution is skewed. \cref{sec:setup} therefore reports the fraction of patients above budget as a separate diagnostic and never presents it as a consequence of \cref{prop:validity}. Second, the slack scales as $L_{\max}\sqrt{\log|\Lambda|/n}$, so a finer grid buys operating-point resolution at a logarithmic cost, and the constant can be sharpened by replacing Hoeffding with an empirical-Bernstein bound, $\sigma(\lambda)\sqrt{2\log(2|\Lambda|/\alpha)/n}+7L_{\max}\log(2|\Lambda|/\alpha)/(3(n-1))$, which is the variant used in all experiments and which reduces to a materially tighter threshold because per-patient alarm counts have low variance. Third, exchangeability across hospitals is an assumption that cross-site transfer violates; the cross-site experiments therefore report the empirical alarm rate rather than asserting that \cref{eq:bound} holds unconditionally.

Two implementation details follow from the proposition. Because $\widehat{R}_n$ need not be monotone on a finite sample, the selector is applied to an isotonic regression of $\widehat{R}_n$ on $\Lambda$; the certificate then refers to the smoothed risk, which is a pool-adjacent weighted average of neighbouring empirical risks rather than a single grid value, so the code exposes this as a switch and reports which variant was used. Second, the bound is vacuous if $L_{\max}\sqrt{\log(|\Lambda|/\alpha)/(2n)}$ approaches $B$, which happens for small calibration sets; the experiments therefore report the ratio of the realised alarm rate to the budget alongside the violation frequency, since a procedure that certifies the budget by emitting almost no alarms is valid but useless. The implementation also clips any loss above $L_{\max}$ and records how often that happened: for an admissible cohort the count is zero, and a non-zero count means $H_{\min}$ was declared too generously, at which point the certificate no longer refers to the quantity that was measured.

\subsection{Operating Point Selection and Lead Time}

Given a calibrated policy, the remaining design question is which feasible threshold to use. Discriminative performance is not the right criterion, because all thresholds satisfying \cref{eq:lambda} already meet the budget; among them the relevant quantity is how early the first alarm arrives relative to onset. For patient $i$ with onset $t^\star_i$, define the lead time
\begin{equation}
\label{eq:lead}
\ell_i(\lambda)=\left(t^{\star}_i-\min\{t:\ a_{it}(\lambda)=1,\ t<t^{\star}_i\}\right)^{+},
\end{equation}
and select $\lambda$ within the feasible set to maximise the median of $\ell_i$ subject to a sensitivity floor, which is the constrained analogue of the usual clinical requirement that a screening rule not miss the majority of events. This is a two-stage procedure: \cref{eq:lambda} certifies feasibility, and \cref{eq:lead} ranks the feasible operating points. \cref{alg:calibration} states the complete procedure.

\begin{algorithm}[tb]
\caption{Blocked conformal calibration of an alarm policy}
\label{alg:calibration}
\begin{algorithmic}[1]
\Require calibration patients $\{(\mathcal{O}_i,t^\star_i,D_i)\}_{i=1}^{n}$; grid $\Lambda$; budget $B$; level $\alpha$; structural parameters $(\tau,\kappa)$
\For{$\lambda\in\Lambda$}
  \State run $\pi_\lambda$ on each calibration patient, obtaining $\{a_{it}\}$
  \State $L_i(\lambda)\leftarrow D_i^{-1}\sum_t a_{it}$
  \State $\widehat{R}_n(\lambda)\leftarrow n^{-1}\sum_i L_i(\lambda)$
\EndFor
\State $\widetilde{R}_n\leftarrow$ isotonic regression of $\widehat{R}_n$ on $\Lambda$ (optional)
\State $\epsilon(\lambda)\leftarrow$ empirical-Bernstein slack of \cref{prop:validity}
\State $\widehat{\lambda}\leftarrow\inf\{\lambda:\widetilde{R}_n(\lambda)+\epsilon(\lambda)\le B\}$
\State \Return $\widehat{\lambda}$, and the empirical violation rate on a held-out split
\end{algorithmic}
\end{algorithm}
